\textbf{Length invariant.} Every finite horizontal/vertical path from S=(0,0) to F=(160,120), using only right and up moves, has horizontal total 160 mm and vertical total 120 mm. Its length is therefore \textbf{280 mm}, regardless of the number or sizes of its pieces. Zero-length pieces do not create turns; adjacent collinear pieces should be merged when counting.

The straight diagonal has length \textbf{200 mm}, since $160^2+120^2=200^2$. Staircases can be made arbitrarily close to it while retaining length 280 mm. Here ``close'' means the largest \textbf{vertical} gap $|y - 3x/4|$ over the whole path, not the length along the route or a perpendicular distance. Closeness of curves does not imply closeness of lengths. No finite staircase becomes the diagonal.

If leftward travel is allowed but downward travel is not, let $L$ be the total distance traveled left. Rightward travel must be $160+L$, so total length is \textbf{280+2L mm}. This counts repeated traversal as traveled length; it is not the area or length of the inked set.

\textbf{Discovery by band.} K-1 compares string lengths and makes close or backtracking routes. Grades 2-3 explains totals and minimum turns. Grades 4-5 separates approximation of location from approximation of length. Measurement suggests equality; regrouping all horizontal and vertical pieces proves it.
